% Work-precision diagrams: rowmap (Rust) vs rowmap.F (Fortran 77).
% Build:  cd docs/report && latexmk -lualatex report.tex
% Figures: uv run plots/work_precision.py --native-only --out docs/report/fig
% Table:   uv run tools/wp_ratio.py --tex docs/report/ratios.tex
\documentclass[11pt,a4paper]{article}
\usepackage{fontspec}
\usepackage{unicode-math}
\usepackage[margin=2.2cm]{geometry}
\usepackage{graphicx}
\usepackage{booktabs}
\usepackage{subcaption}
\usepackage{hyperref}

\title{ROWMAP in Rust: work-precision diagrams}
\author{rowmap-ng}
\date{31 October 2026}

\newcommand{\fig}[2]{%
  \begin{subfigure}{0.49\textwidth}\centering
    \includegraphics[width=\linewidth]{fig/wp_#1.pdf}\caption{#2}
  \end{subfigure}}

\begin{document}
\maketitle

\section{Setup}

All runs use the paper grids, $\mathrm{rtol}=\mathrm{atol}=\mathrm{tol}$ for
$\mathrm{tol}=10^{-2},\dots,10^{-7}$, and a maximum Krylov dimension of~66.
$\mathrm{ERR}$ is the paper's error measure at $t_\mathrm{end}$ against a
DOPRI5 reference ($\mathrm{tol}=10^{-10}$). The CPU time covers the
integration only. It is the mean over repeated runs that take at least
0.5\,s together (Fortran driver and Rust example alike). The machine is
an Intel Core i5-6500 (AVX2/FMA) with Rust~1.98 (\texttt{--release}, no
\texttt{target-cpu}; the AVX2/FMA kernels are chosen at run time), and
gfortran~16.2 \texttt{-O2} with OpenBLAS.

Series:
\begin{description}
  \item[F77 GRK4T] the original \texttt{rowmap.F}, exact $J v$; the
    non-autonomous problems use $f_t$ (strategy sft${}=2$).
  \item[Rust GRK4T] default method, MAP schedule with the minimal dimensions
    4/7/8/11 (autonomous) and 4/9/10/15 (non-autonomous).
  \item[Rust ROS34PW2] the order-3 option (\texttt{Method::Ros34pw2},
    Rang--Angermann 2005: stiffly accurate, L-stable), minimal dimensions
    3/4/4/4 and 3/4/8/8.
  \item[Rust GRK4T (no-stage4-closure)] diagnostic only, autonomous problems:
    the schedule without the stage-4 closure (4/7/7/8). It keeps classical
    order~4 (\texttt{docs/map.md} §4); this is open question~2 for the
    authors.
\end{description}

\section{Work-precision diagrams}

\begin{figure}[htbp]\centering
  \fig{heat_m50}{HEAT}\hfill\fig{nilidi_m69}{NILIDI}\\[1ex]
  \fig{bruss_m50}{BRUSS}\hfill\fig{food_m50}{FOOD}\\[1ex]
  \fig{vdpl_m20}{VDPL}\hfill
  \begin{subfigure}{0.49\textwidth}\end{subfigure}
  \caption{Autonomous problems: $\lg\mathrm{ERR}$ versus CPU time.}
  \label{fig:aut}
\end{figure}

\begin{figure}[htbp]\centering
  \fig{difkon_m69}{DIFKON}\hfill\fig{diffu2_m69}{DIFFU2}
  \caption{Non-autonomous problems (native form, exact $f_t$).}
  \label{fig:nonaut}
\end{figure}

\section{Summary}

Table~\ref{tab:ratio} compares CPU time at \emph{equal accuracy}: for each
Fortran point, the Rust CPU time at the same ERR is interpolated (log-log)
from the Rust curve. The number in parentheses counts the Fortran points
inside the Rust ERR range.

\begin{table}[htbp]\centering
  \begin{tabular}{lrrr}
\toprule
Problem & GRK4T & GRK4T reduced & ROS34PW2 \\
\midrule
HEAT & 1.19 (5) & 0.97 (5) & 0.57 (5) \\
NILIDI & 1.52 (3) & 1.44 (3) & 1.93 (5) \\
BRUSS & 0.92 (5) & 0.74 (6) & 1.01 (6) \\
FOOD & 1.36 (6) & 0.91 (6) & 1.16 (6) \\
VDPL & 1.56 (5) & 1.11 (5) & 1.94 (3) \\
DIFKON & 1.35 (5) & -- & 2.46 (5) \\
DIFFU2 & 1.21 (6) & -- & 0.90 (2) \\
\bottomrule
\end{tabular}

  \caption{CPU time Rust/Fortran at equal ERR (geometric mean). Values below
    1 mean Rust is faster.}
  \label{tab:ratio}
\end{table}

\begin{itemize}
  \item \textbf{GRK4T, Rust vs Fortran.} Same step counts (within a few
    percent; NILIDI and the non-autonomous problems need 15--25\,\% fewer
    steps). ERR at the same tolerance is equal or smaller. At equal accuracy
    Rust takes 0.92--1.56 times the Fortran CPU time; it is faster on BRUSS.
  \item \textbf{ROS34PW2} is the method of choice at crude tolerances
    (HEAT: 0.57, DIFFU2: 0.90). At tighter tolerances its lower order costs
    steps; on NILIDI, VDPL and DIFKON it is about twice as expensive as GRK4T.
  \item \textbf{The reduced schedule} saves 5--30\,\% CPU at equal accuracy
    and beats Fortran on HEAT, BRUSS and FOOD.
\end{itemize}

\section{The stage systems: incremental Givens QR}

$H=Q^TAQ$ has a generalized Hessenberg structure (\texttt{docs/map.md} §9,
proved in Lean, \texttt{HStructure.lean}). A closed column is final and zero
below its closing row, and the closed columns form a prefix. All stages of a
step share $M=I-\gamma hH$. So each closed column is rotated to triangular
form exactly once. A solve treats only the few open trailing columns, and the
stage-1 residual test needs only their solution components, with no full back
substitution. This replaces the former dense LU from scratch per stage and
per Arnoldi iteration. On VDPL the share of the small systems drops from
11\,\% to 3.7\,\% of the run time; elsewhere it was small anyway.

\section{Compiler and memory: the same algorithm}

\texttt{rowmap.F} and the Rust code differ in the algorithm. Rust completes
the proven MAP schedule and re-orthogonalizes; \texttt{rowmap.F} stops
stages 2--4 early and does a single MGS pass. To compare compiler and memory
behaviour, Rust was also run with a \emph{single} classical Gram--Schmidt pass
(experiment flag \texttt{--single-pass}). The CPU time was normalized by the
number of $Jv$ products, since each $Jv$ comes with one orthogonalization
against the current basis (Table~\ref{tab:jv}).

\begin{table}[htbp]\centering\small
  \begin{tabular}{lrrrrrrrrr}
\toprule
& \multicolumn{3}{c}{F77} & \multicolumn{3}{c}{Rust 1-pass} & \multicolumn{3}{c}{Rust 2-pass} \\
Problem & $\mu$s/$Jv$ & $Jv$/step & ratio & $\mu$s/$Jv$ & $Jv$/step & ratio & $\mu$s/$Jv$ & $Jv$/step & ratio \\
\midrule
HEAT & 34.5 & 18.5 & 1.00 & 27.2 & 21.8 & 0.79 & 38.8 & 19.5 & 1.13 \\
NILIDI & 118.9 & 42.3 & 1.00 & 129.4 & 51.8 & 1.09 & 181.5 & 49.8 & 1.53 \\
BRUSS & 67.6 & 8.4 & 1.00 & 50.6 & 12.8 & 0.75 & 63.0 & 11.5 & 0.93 \\
FOOD & 66.3 & 10.3 & 1.00 & 51.8 & 15.9 & 0.78 & 69.7 & 13.7 & 1.05 \\
VDPL & 11.9 & 15.3 & 1.00 & 10.7 & 21.1 & 0.90 & 15.9 & 17.7 & 1.34 \\
DIFKON & 104.1 & 31.9 & 1.00 & 99.7 & 40.3 & 0.96 & 136.0 & 40.2 & 1.31 \\
DIFFU2 & 106.2 & 56.5 & 1.00 & 100.4 & 63.3 & 0.95 & 156.7 & 63.2 & 1.48 \\
\bottomrule
\end{tabular}

  \caption{Cost per $Jv$ product (user $Jv$ + Gram--Schmidt + bookkeeping),
    geometric means over $\mathrm{tol}=10^{-2},\dots,10^{-7}$. Ratio:
    Rust/Fortran per $Jv$.}
  \label{tab:jv}
\end{table}

\begin{itemize}
  \item \textbf{With a single pass, Rust is faster per $Jv$ on six of seven
    problems} (ratio 0.75--0.96; NILIDI 1.09). This holds although Rust builds
    12--54\,\% \emph{more} vectors per step, so each orthogonalization runs
    against a larger basis. The kernels are blocked, use AVX2/FMA, and keep
    $Q$ contiguous; \texttt{rowmap.F} calls BLAS-1 (\texttt{ddot},
    \texttt{daxpy}) once per basis vector. NILIDI, with the longest basis
    ($\kappa\approx50$), is the exception; the cause was not investigated.
  \item \textbf{The second pass costs 25--56\,\% per $Jv$} (two-pass ratio
    0.93--1.53). Together with the 5--37\,\% more $Jv$ per step (complete
    schedule), this accounts for the gap in Table~\ref{tab:ratio}.
  \item \textbf{Single pass is not an option for the solver.} With one pass,
    13 unit tests fail. At (near) breakdown the normalized rounding remainder
    is not orthogonal to $Q$ ($\|Q^TQ-I\|=1$, step results off by $O(10)$),
    and $w_4\in K_3$ is no longer detected. On the paper problems, VDPL
    loses a factor of 30 in ERR at $\mathrm{tol}=10^{-7}$, and HEAT needs 254
    instead of 99 steps. The second pass stays the default (decision of the
    authors); the single pass remains in the code for comparisons only.
\end{itemize}

\section{Room for optimization}

Profile (feature \texttt{profile}, $\mathrm{tol}=10^{-5}$): orthogonalization
takes 37--57\,\% of the run time, the user's $Jv$ 16--31\,\%, $f$ 1--7\,\%,
the Givens QR at most 4\,\% (VDPL), and the rest a few percent. The kernels run near the
memory and FMA limits. So what remains is algorithmic, and each option needs
a decision:

\begin{enumerate}
  \item \textbf{Stage-4 closure} (GRK4T autonomous). Dropping it is
    classically sufficient for order~4 and saves 5--30\,\%. It is the only
    option that keeps an order guarantee (the classical one).
  \item \textbf{Early stop in stages 2--4} as in \texttt{rowmap.F}: fewer
    $Jv$, but no longer the proven minimal dimensions.
  \item \textbf{Smaller items}, a few percent each: allocation reuse,
    \texttt{target-cpu=native}.
\end{enumerate}

\end{document}
